% =============================================================================
\chapter{The Mesh Library}
\label{ch:mesh}
% =============================================================================

\section{Overview}
\label{sec:mesh:overview}

% Role of the mesh library, relationship to FE library

\section{Mesh data structures}
\label{sec:mesh:data}

% Vertices, cells, faces, edges; connectivity arrays

\section{Element types}
\label{sec:mesh:elements}

% Triangles, quadrilaterals, tetrahedra, hexahedra, wedges, pyramids

\section{Mesh input and output}
\label{sec:mesh:io}

% VTK/VTU format, mesh readers

\section{Mesh topology and adjacency}
\label{sec:mesh:topology}

% Cell-to-vertex, face-to-cell, graph coloring for parallel assembly

\section{Mesh partitioning for parallel computation}
\label{sec:mesh:partitioning}

% Domain decomposition, ghost layers, MPI distribution

\section{Geometric queries}
\label{sec:mesh:queries}

% Point location, bounding boxes, spatial search
